%%%%%%%%%%%%%%%%%%%%%%%%%%%%%%%%%%%%%%%%%%%%%%%%%%%%%%%%%%%%%%%%%%%
\section{Mapping the Detector into a CAM}
\label{sec:Mapping}
%%%%%%%%%%%%%%%%%%%%%%%%%%%%%%%%%%%%%%%%%%%%%%%%%%%%%%%%%%%%%%%%%%%
% Restructured Sep 1 per MN's comment (D): the "What CAM-resident
% requires" block now LEADS the section (one move, per MN), and the
% encodings are presented as options for satisfying those
% requirements. Order: intro (reframed per MN: technology-enabled
% first, CMOS as the nearer-term study) -> requirements -> first
% mapping -> path back -> encodings 1-3 -> decision rule.
% Aug 31 structure otherwise unchanged: mappings only here, all
% evaluation in Sec. V.

This section covers ways to map activation signatures onto \cam
arrays. Candidate arrays span emerging devices through nearer-term
CMOS solutions that leverage different encoding schemes -- notably,
traditional \cam designs can also be realized with nonvolatile
memories (RRAM, MRAM, etc.), which could approach the density gains
and, with a mostly static signature bank, benefit from
nonvolatility. Sec.~\ref{subsec:Requirements} first spells out what
any \cam-resident realization requires. The first mapping
(Sec.~\ref{subsec:FirstPipeline}) assumed the best-case device -- a
multi-bit, \fefet-class cell whose matchline computes the Euclidean
distance the full-precision detector actually uses -- and it is
retained as the reference point the other encodings are measured
against. The encodings of
Secs.~\ref{subsec:SignRP}--\ref{subsec:TwoStage} satisfy the
requirements on hardware buildable today, including Hamming-distance
options that are near-term and realizable by a host of NVMs as well
as by plain CMOS.
\inote[aiviolet]{Sep 2, your (B): opener rewritten generically
(``ways to map activation signatures to CAM arrays''), your
right-margin NVM wording worked in, reference-point sentiment kept.
The old ``two constraints'' sentence (deployment target + audit) is
cut here since IV-A's closing paragraph already carries the audit
point -- say the word if you want it restored}

\subsection{What ``CAM-resident'' requires}
\label{subsec:Requirements}
For the monitor to be one memory structure rather than a memory plus an
accelerator, three properties have to hold. {\bf (1)} The query has to
be producible from the raw activation by something cheap, since whatever
sits in front of the array runs on every token. {\bf (2)} The comparison
the array performs natively has to be the comparison the detector needs,
which for a binary or ternary cell means the decision has to be
expressible as a Hamming distance. {\bf (3)} The decision rule has to
consume what a matchline can report -- a winner, a set of rows under a
threshold, or a count -- rather than a full sorted score
vector.\footnote{Returning all $K$ scores is a readout job a \cam is
not for.}

Property~(1) also has a methodological consequence. If the front end is
trained, it can learn to do the classification itself, leaving the
memory as decoration. Sec.~\ref{subsec:Audit} describes the test we
use to detect that, and what it caught when we ran it on pipelines
that did have trained front ends.

\subsection{The first mapping: learned projection, 3-bit Euclidean}
\label{subsec:FirstPipeline}
A multi-bit, \fefet-class cell broadens what a stored row can be:
each cell holds a level rather than a bit, which could allow for a
richer representation of the data -- or the same representation in
fewer memory cells -- with the matchline computing the Euclidean
distance the detector actually uses (Sec.~\ref{subsec:CAMspace}). A
down-projection mechanism does need to be implemented regardless,
since whatever sits in front of the array runs on every token
(requirement~{\bf (1)}): each raw activation is reduced from $d$
dimensions to a much smaller $d'$ and quantized to a few bits per
dimension.\inote[aiviolet]{Sep 2: opening reframed per your margin
note -- leads with what \fefet \cam{}s can do (richer representation
or fewer cells), with the down-projection-regardless point kept}
Concretely, the pipeline took the residual-stream
activation at layer 33 (the tuned layer choice described in
Sec.~\ref{subsec:Setup}), built a bank of mutually distinct deception
directions by iterative nullspace
projection~\cite{ravfogel2020inlp} (train a probe, project its direction
out, retrain), weighted each stored direction with one trained scalar,
compressed the activation with a \emph{learned} down-projection from
$d=8{,}192$ to $d'=128$ dimensions, and quantized query and stored rows
alike to 3 bits with learned-step quantization~\cite{esser2020lsq}, so
that the low-precision rounding sits inside the training loop rather
than being applied afterwards.\footnote{Quantization-aware training at
this precision is standard practice: a shipping on-device foundation
model is trained with 2-bit QAT~\cite{apple2025foundation}.} Matching
was Euclidean winner-take-all, returning the winning row and its
margin. $d'=128$ was chosen to match the row geometry of a multi-bit
\fefet array -- 128 cells sharing a matchline, at 3 bits per cell --
so the stored row is a 384-bit-\emph{equivalent} row (physically, 128
cells), which is where the matched 384-bit operating point used
throughout Sec.~\ref{sec:Results} comes
from.\inote[aiviolet]{Sep 1, three of your notes on this passage:
(1) it now gives the strategy first ($d \to d'$, then quantize) and
the numbers after; (2) ``384-bit row'' is now ``384-bit-equivalent
row,'' per your ``there are just 128 cells''; (3) the why-128
sentence ties $d'$ to \fefet row geometry per your margin note --
the 128-cells-per-matchline reading is mine, so fix the parenthetical
if the real constraint is different. Layer 33's proper introduction
moved to the tutorial setup of Sec.~V-A; this passage now just points
there}

This mapping's results -- and the audit finding that retired its
trained cousins -- are evaluated in Sec.~\ref{sec:Results}
(Fig.~\ref{fig:grids}(b) and
Table~\ref{tab:audit}).\inote[aiviolet]{Sep 1: the rest of this
paragraph (``...the encodings in the rest of this section descend
from the variants whose stored rows carry the classification'') is
cut per your ``can probably cut/scale this''}

\label{subsec:PathBack}%
The encodings that follow trade this mapping's multi-bit cell for
Hamming-distance solutions, which are near-term and realizable by a
host of NVMs as well as by plain CMOS; the natural recombination -- a
frozen projection matched on a multi-bit Euclidean matchline -- has
now been run on the deception grid, and with the projection frozen the
Euclidean matchline buys no accuracy over Hamming distance
(Sec.~\ref{subsec:BitWidth}).\ib{Item 8(a) run; result in Sec.~V-C.}
\inote[aiviolet]{Sep 2: the old Sec.~IV-C (``The path back to a
multi-bit cell'') is cut to this transition sentence per your notes
-- the ``abandonment'' framing you flagged as making no sense is
gone with it, and the Hamming/NVM near-term point you asked for is
in. The full old text is preserved as a comment in the source; the
recombination-experiment question stays in the review appendix, and
Sec.~V's audit subsection still points here}
% --- Old Sec. IV-C, cut Sep 2 per MN (kept for possible restore) ---
% Moving to frozen projections and a Hamming metric is not an
% abandonment of the Euclidean track, and the two can be recombined.
% The property that made the first pipeline suspect was the learned
% front end, not the multi-bit cell; the property that makes the new
% encodings trustworthy is the frozen projection, not the Hamming
% metric. A frozen random projection quantized to 3 bits and matched
% on a multi-bit Euclidean matchline would have both: nothing trained
% ahead of the array, and a metric that keeps magnitude rather than
% only the sign of each coordinate. It is worth trying because sign
% codes discard exactly the information the close calls depend on,
% and because if magnitude carries signal, a few hundred 3-bit values
% may replace 2,048 sign bits, i.e., a smaller bank for the same
% separation. / Two reasons it might not work are worth stating in
% advance. A learned projection can be trained through the quantizer,
% whereas a frozen one cannot, so the 3-bit levels have to be placed
% after the fact and the frozen projection may not survive
% quantization as gracefully. And it gives up the argument that the
% whole detector fits in a plain CMOS array today.
% -----------------------------------------------------------------

\subsection{Encoding 1: sign random projection, ternary (or via other
NVMs)}
\label{subsec:SignRP}
A first option for meeting requirements {\bf (1)}--{\bf (3)} on
hardware buildable today is a fixed random projection: draw
$P \in \mathbb{R}^{m \times d}$ once and freeze it, then encode an
activation as the sign pattern of its projection,
%
\begin{equation}
c_j(h) \;=\; \operatorname{sign}\!\big( (Ph)_j \big), \qquad j = 1 \ldots m,
\label{eq:signrp}
\end{equation}
%
with a dead zone $|(Ph)_j| < \epsilon$ mapped to don't-care, which makes
the code ternary and requires a wildcard-capable cell. This is the
standard sign-random-projection (SimHash)
construction~\cite{charikar2002simhash}; the Hamming distance between
two codes estimates the angle between the underlying
vectors.\inote[aiviolet]{Sep 2, per your marks: opener simplified
(``too AI-like''), the hyperplane explanation shortened per ``CAM
simplify,'' and the ``contribution here is not the construction''
sentence deleted per your strike}

From a hardware perspective, the comparison is an XOR plus a
population count, which a plain
\cam array computes natively. A code wider than a practical array row
is split across subarrays and the results aggregated, which is
standard practice and a real design parameter -- the measured
\fefet \cam demonstration of~\cite{kazemi2022hdcfefet} uses exactly
this organization, with voting to aggregate results between
subarrays.

\subsection{Encoding 2: thermometer coding (plain binary CMOS)}
\label{subsec:Thermometer}
A second option for reaching a \cam realization -- and one a plain
CMOS binary array can serve: quantize each projected dimension into
$\ell$ levels and encode each
level in unary, so a value of level $v$ becomes $v$ ones followed by
$\ell-v$ zeros. For two codes $u,v$ over the same dimension,
%
\begin{equation}
d_{\mathrm{Ham}}\big(\mathrm{therm}(u), \mathrm{therm}(v)\big)
\;=\; \sum_i \big| u_i - v_i \big|,
\label{eq:thermo}
\end{equation}
%
i.e., Hamming distance on the code is exactly the $L_1$ distance on the
quantized values. The identity is exact, not approximate, which is what
makes the encoding useful here: a cell that can only report disagreement
now supports a magnitude-sensitive metric. No wildcard is needed, so
the cell is a plain binary \cam at roughly 6--10 transistors rather than
a \tcam at roughly 16.

The cell menu is wider than CMOS. A 1\fefet-1C array -- a DRAM-like
cell with the \fefet as access transistor -- has been extended to
ternary \cam operation, with a fabricated proof of concept whose
sensed matchline voltage after charge sharing is linear in Hamming
distance~\cite{fefet1c2026}. A plausible extension is to host either
of this section's encodings on such arrays: the Hamming linearity
reads off the distances of Eqs.~(\ref{eq:signrp})
and~(\ref{eq:thermo}) directly, the don't-care support serves the
ternary code's dead zone, and the DRAM-like cell promises density --
especially considering rising 3D DRAM -- with nonvolatile storage
suited to a mostly static signature bank.\inote[aiviolet]{Sep 2,
your two-arrow note at the identity claim: the 1FeFET-1C material is
in, from the TED draft you attached (cited as ``manuscript in
preparation'' -- fix the entry when the author list settles). The
DAC-thermosensitive.pdf citation and its extension still need the
file from work}
\dt[DAC-thermosensitive.pdf: cite here with its plausible extension
once the PDF arrives from MN's work machine; also settle the
fefet1c2026 bib entry's authors before submission]{}

The comparison against \signrp is therefore not symmetric. At a matched
code length the two encodings use the same number of \emph{bits} but not
the same amount of \emph{silicon}, and Sec.~\ref{sec:Results} reports
both alongside that qualification.

\subsection{The decision rule}

\label{subsec:Decision}
The stored rows are labeled example activations rather than fitted
directions -- i.e., each domain stores not one direction per class
but a few thousand encoded activations, each labeled honest or
deceptive. A query is scored by the fraction of its $k$ nearest stored
rows, under the array's metric, that carry the deceptive label, with
$k=31$ in the runs reported here (the gain flattens well before 31);
\auroc is computed on that score directly, so no threshold is
chosen.\ib{Item 7: this is what the runs do. Each domain's 70\% split
goes through the frozen encoder (mean-center, random projection to 128
dims, quantize to 4 levels, thermometer code) and is stored with its
label -- 831 rows for definitional up to 3{,}201 for sycophancy. A test
activation is encoded the same way, compared by Hamming distance
against every stored row, and scored by the fraction of its 31 nearest
rows labeled deceptive; AUROC is computed on that fraction, so there is
no majority vote and no threshold. The bank description, the audit
discussion, and Table II's row counts were already right; only
``majority label'' was imprecise.} The scheme is analogous to few-shot classification:
rather than one fitted direction per behavior (the $N$-way, 1-shot
analog), the bank operates $N$-way, many-shot, with thousands of
stored exemplars per domain and the array's distance metric supplying
the comparison. A vote over neighbors rather than a single winner is
worth 3--5 accuracy points in in-distribution tests, and it consumes
only what a best-match or threshold matchline can report.

\subsection{Encoding 3, as a ceiling: shortlist plus exact rerank}
\label{subsec:TwoStage}
To bound what the compression costs rather than only measuring it, we
also run a two-stage variant: stage 1 is a thermometer-Hamming search of
the whole bank that keeps the top $k$ rows as a shortlist, and stage 2
is an exact cosine rerank on the raw, uncompressed activations
restricted to that shortlist. With $k$ equal to the whole bank, stage 2
alone defines a full-precision ceiling for the same data and protocol.
Shortlist-then-rerank is standard practice in large-scale retrieval
(e.g., the two-stage GPU indexes of~\cite{faisscuvs2025});
equivalently, stage 2 can be read as checking a smaller $K$ that a
GPU can easily and accurately serve, with stage 1 reducing the
library to that size. The variant is also the architecture a real
system might adopt if
single-stage fidelity is inadequate, at the cost of a second memory
structure and a floating-point stage.
